\documentclass[10pt,a4paper]{amsart}
\usepackage{amsmath,amssymb,amsthm,graphicx}
% Retained preamble for the Stacks Project chapter snapshot in this directory.
% The chapter's own preamble.tex is a copy of the shared Stacks Project preamble; the block
% below is sliced MECHANICALLY, line for line, from it. Nothing was added, removed or
% redefined; the wrapper class, the xr-hyper/hyperref wiring and the externaldocument list
% are omitted because the assembler supplies the document class.
%
% lines taken from preamble.tex: 3, 5, 9, 11, 12, 156, 157, 158, ... (51 lines total)
%
% package coverage restored from the chapter's own preamble (2026-10-13): xy, xy-2cell

\usepackage[all]{xy}
\xyoption{2cell}
\UseAllTwocells

\newenvironment{reference}{\comment}{\endcomment}
\newenvironment{slogan}{\comment}{\endcomment}
\newenvironment{history}{\comment}{\endcomment}
\theoremstyle{plain}
\newtheorem{theorem}[subsection]{Theorem}
\newtheorem{proposition}[subsection]{Proposition}
\newtheorem{lemma}[subsection]{Lemma}
\theoremstyle{definition}
\newtheorem{definition}[subsection]{Definition}
\newtheorem{example}[subsection]{Example}
\newtheorem{exercise}[subsection]{Exercise}
\newtheorem{situation}[subsection]{Situation}
\theoremstyle{remark}
\newtheorem{remark}[subsection]{Remark}
\newtheorem{remarks}[subsection]{Remarks}
\numberwithin{equation}{subsection}
\def\lim{\mathop{\mathrm{lim}}\nolimits}
\def\colim{\mathop{\mathrm{colim}}\nolimits}
\def\Spec{\mathop{\mathrm{Spec}}}
\def\Hom{\mathop{\mathrm{Hom}}\nolimits}
\def\Ext{\mathop{\mathrm{Ext}}\nolimits}
\def\SheafHom{\mathop{\mathcal{H}\!\mathit{om}}\nolimits}
\def\SheafExt{\mathop{\mathcal{E}\!\mathit{xt}}\nolimits}
\def\Sch{\mathit{Sch}}
\def\Mor{\mathop{\mathrm{Mor}}\nolimits}
\def\Ob{\mathop{\mathrm{Ob}}\nolimits}
\def\Sh{\mathop{\mathit{Sh}}\nolimits}
\def\NL{\mathop{N\!L}\nolimits}
\def\CH{\mathop{\mathrm{CH}}\nolimits}
\def\proetale{{pro\text{-}\acute{e}tale}}
\def\etale{{\acute{e}tale}}
\def\QCoh{\mathit{QCoh}}
\def\Ker{\mathop{\mathrm{Ker}}}
\def\Im{\mathop{\mathrm{Im}}}
\def\Coker{\mathop{\mathrm{Coker}}}
\def\Coim{\mathop{\mathrm{Coim}}}
\DeclareMathSymbol{\boxtimes}{\mathbin}{AMSa}{"02}
\def\QCohstack{\mathcal{QC}\!\mathit{oh}}
\def\Cohstack{\mathcal{C}\!\mathit{oh}}
\def\Spacesstack{\mathcal{S}\!\mathit{paces}}
\def\Quotfunctor{\mathrm{Quot}}
\def\Hilbfunctor{\mathrm{Hilb}}
\def\Curvesstack{\mathcal{C}\!\mathit{urves}}
\def\Polarizedstack{\mathcal{P}\!\mathit{olarized}}
\def\Complexesstack{\mathcal{C}\!\mathit{omplexes}}
\def\Pic{\mathop{\mathrm{Pic}}\nolimits}
\def\Picardstack{\mathcal{P}\!\mathit{ic}}
\def\Picardfunctor{\mathrm{Pic}}
\def\Deformationcategory{\mathcal{D}\!\mathit{ef}}
\usepackage{cleveref}
\crefname{originallemma}{original lemma}{original lemmas}
\Crefname{originallemma}{Original lemma}{Original lemmas}
\makeatletter
\newlabel{algebra-lemma-cover-module}{{00EK}{source}{}{sourcepointer-1}{}}
\newlabel{algebra-lemma-cover-module@cref}{{[originallemma][0][]00EK}{[1][0][]source}{}{}{}}
\makeatother
\crefname{originallemma}{original lemma}{original lemmas}
\Crefname{originallemma}{Original lemma}{Original lemmas}
\makeatletter
\newlabel{lemma-glue-trace}{{0BTA}{source}{}{sourcepointer-2}{}}
\newlabel{lemma-glue-trace@cref}{{[originallemma][0][]0BTA}{[1][0][]source}{}{}{}}
\makeatother
\crefname{originallemma}{original lemma}{original lemmas}
\Crefname{originallemma}{Original lemma}{Original lemmas}
\makeatletter
\newlabel{lemma-trace-base-change}{{0BT9}{source}{}{sourcepointer-3}{}}
\newlabel{lemma-trace-base-change@cref}{{[originallemma][0][]0BT9}{[1][0][]source}{}{}{}}
\makeatother
\crefname{originallemma}{original lemma}{original lemmas}
\Crefname{originallemma}{Original lemma}{Original lemmas}
\makeatletter
\newlabel{lemma-trace-unique}{{0BT7}{source}{}{sourcepointer-4}{}}
\newlabel{lemma-trace-unique@cref}{{[originallemma][0][]0BT7}{[1][0][]source}{}{}{}}
\makeatother
\title{Stacks, Discriminants: lemma-dualizing-tau}
\author{The Stacks Project Authors}
\date{stacks-project:discriminant.tex:local-snapshot-20261012; GNU FDL 1.2 or later}
\begin{document}
\maketitle
\noindent\textit{Editorial note.} Statement span is lines 997-1001 of sources/stacks-discriminant/pinned-original-source.tex; proof span is lines 1003-1104, the complete author proof (100 body lines).\par\medskip
\noindent\textit{Editorial note.} Navigation stubs: 4 cited-but-not-included labels are supplied as non-load-bearing pointers with their pinned tags where the snapshot tag index has them.\par\medskip
\noindent\textit{Editorial note.} extra\_packages is empty on purpose: the retained preamble already loads xy with options and a second declaration clashes in the assembler wrapper.\par\medskip
\noindent\textit{Editorial note.} Text is retained verbatim, including the author's own typos and placeholders.\par\medskip
\noindent\textit{Editorial note.} Queue index 1 of queue\_discriminant.json (pb=4592, stmt 997-1001, proof 1003-1104).\par\medskip
\noindent\textit{Allowed standard prerequisites (editorial).}\begin{itemize}
\item Explicit prerequisite boundary. The cited inputs are prerequisites; the argument in the retained span is not replaced by any cited result.
\end{itemize}\medskip
\ifdefined\hypertarget\hypertarget{sourcepointer-1}{}\fi
\noindent\textit{Source reference.} Original lemma 00EK: sources/stacks-discriminant/pinned-original-source.tex, source lines 1104--1104. Cited by the retained statement or proof but not part of the collected spans; supplied only as a navigation pointer. The argument does not depend on its content.\par\medskip
\ifdefined\hypertarget\hypertarget{sourcepointer-2}{}\fi
\noindent\textit{Source reference.} Original lemma 0BTA: sources/stacks-discriminant/pinned-original-source.tex, source lines 955--961. Cited by the retained statement or proof but not part of the collected spans; supplied only as a navigation pointer. The argument does not depend on its content.\par\medskip
\ifdefined\hypertarget\hypertarget{sourcepointer-3}{}\fi
\noindent\textit{Source reference.} Original lemma 0BT9: sources/stacks-discriminant/pinned-original-source.tex, source lines 894--907. Cited by the retained statement or proof but not part of the collected spans; supplied only as a navigation pointer. The argument does not depend on its content.\par\medskip
\ifdefined\hypertarget\hypertarget{sourcepointer-4}{}\fi
\noindent\textit{Source reference.} Original lemma 0BT7: sources/stacks-discriminant/pinned-original-source.tex, source lines 843--846. Cited by the retained statement or proof but not part of the collected spans; supplied only as a navigation pointer. The argument does not depend on its content.\par\medskip
\begin{lemma}
\label{lemma-dualizing-tau}
Let $A \to B$ be a flat quasi-finite map of Noetherian rings.
There exists a trace element $\tau \in \omega_{B/A}$.
\end{lemma}

\begin{proof}
Choose a factorization $A \to B' \to B$ with $A \to B'$ finite and
$\Spec(B) \to \Spec(B')$ an open immersion. Let $g_1, \ldots, g_n \in B'$
be elements such that $\Spec(B) = \bigcup D(g_i)$ as opens of $\Spec(B')$.
Suppose that we can prove the existence of trace elements $\tau_i$ for the
quasi-finite flat ring maps $A \to B_{g_i}$. Then for all $i, j$ the elements
$\tau_i$ and $\tau_j$ map to trace elements of $\omega_{B_{g_ig_j}/A}$
by Lemma \ref{lemma-trace-base-change}. By uniqueness of
trace elements (Lemma \ref{lemma-trace-unique}) they map to the same element.
Hence the sheaf condition for the quasi-coherent module associated to
$\omega_{B/A}$ (see Algebra, Lemma \ref{algebra-lemma-cover-module})
produces an element $\tau \in \omega_{B/A}$.
Then $\tau$ is a trace element by
Lemma \ref{lemma-glue-trace}.
In this way we reduce to the case treated in the next paragraph.

\medskip\noindent
Assume we have $A \to B'$ finite and $g \in B'$ with $B = B'_g$ flat over $A$.
It is our task to construct a trace element in
$\omega_{B/A} = \Hom_A(B', A) \otimes_{B'} B$.
Choose a resolution $F_1 \to F_0 \to B' \to 0$ of $B'$ by finite free
$A$-modules $F_0$ and $F_1$. Then we have an exact sequence
$$
0 \to \Hom_A(B', A) \to F_0^\vee \to F_1^\vee
$$
where $F_i^\vee = \Hom_A(F_i, A)$ is the dual finite free module.
Similarly we have the exact sequence
$$
0 \to \Hom_A(B', B') \to F_0^\vee \otimes_A B' \to F_1^\vee \otimes_A B'
$$
The idea of the construction of $\tau$ is to use the diagram
$$
B' \xrightarrow{\mu} \Hom_A(B', B')
\leftarrow \Hom_A(B', A) \otimes_A B'
\xrightarrow{ev} A
$$
where the first arrow sends $b' \in B'$ to the $A$-linear operator
given by multiplication by $b'$ and the last arrow is the evaluation map.
The problem is that the middle arrow, which sends $\lambda' \otimes b'$
to the map $b'' \mapsto \lambda'(b'')b'$, is not an isomorphism.
If $B'$ is flat over $A$, the exact sequences above show that it
is an isomorphism and the composition from left to right is the usual trace
$\text{Trace}_{B'/A}$. In the general case, we consider
the diagram
$$
\xymatrix{
& \Hom_A(B', A) \otimes_A B' \ar[r] \ar[d] &
\Hom_A(B', A) \otimes_A B'_g \ar[d] \\
B' \ar[r]_-\mu \ar@{..>}[rru] \ar@{..>}[ru]^\psi &
\Hom_A(B', B') \ar[r] &
\Ker(F_0^\vee \otimes_A B'_g \to F_1^\vee \otimes_A B'_g)
}
$$
By flatness of $A \to B'_g$ we see that the right vertical arrow is an
isomorphism. Hence we obtain the unadorned dotted arrow.
Since $B'_g = \colim \frac{1}{g^n}B'$, since
colimits commute with tensor products,
and since $B'$ is a finitely presented $A$-module
we can find an $n \geq 0$ and a $B'$-linear (for right $B'$-module structure)
map $\psi : B' \to \Hom_A(B', A) \otimes_A B'$
whose composition with the left vertical arrow is $g^n\mu$.
Composing with $ev$ we obtain an element
$ev \circ \psi \in \Hom_A(B', A)$. Then we set 
$$
\tau = (ev \circ \psi) \otimes g^{-n} \in
\Hom_A(B', A) \otimes_{B'} B'_g = \omega_{B'_g/A} = \omega_{B/A}
$$
We omit the easy verification that this element does not depend
on the choice of $n$ and $\psi$ above.

\medskip\noindent
Let us prove that $\tau$ as constructed in the previous paragraph
has the desired property in a special case. Namely, say
$B' = C' \times D'$ and $g = (f, h)$ where $A \to C'$ flat, $D'_h$ is flat, and
$f$ is a unit in $C'$.
To show: $\tau$ maps to $\text{Trace}_{C'/A}$ in $\omega_{C'/A}$.
In this case we first choose $n_D$ and
$\psi_D : D' \to \Hom_A(D', A) \otimes_A D'$ as above for the pair
$(D', h)$ and we can let
$\psi_C : C' \to \Hom_A(C', A) \otimes_A C' = \Hom_A(C', C')$
be the map seconding $c' \in C'$ to multiplication by $c'$.
Then we take $n = n_D$ and $\psi = (f^{n_D} \psi_C, \psi_D)$
and the desired compatibility is clear because
$\text{Trace}_{C'/A} = ev \circ \psi_C$ as remarked above.

\medskip\noindent
To prove the desired property in general, suppose given
$A \to A_1$ with $A_1$ Noetherian and a product decomposition
$B'_g \otimes_A A_1 = C \times D$ with $A_1 \to C$ finite.
Set $B'_1 = B' \otimes_A A_1$. Then $\Spec(C) \to \Spec(B'_1)$
is an open immersion as $B'_g \otimes_A A_1 = (B'_1)_g$ and
the image is closed as $B'_1 \to C$ is finite (since $A_1 \to C$
is finite). Thus $B'_1 = C \times D'$ and $D'_g = D$.
We conclude that $B'_1 = C \times D'$ and $g$ over $A_1$
are as in the previous paragraph.
Since formation of the displayed diagram above
commutes with base change, the formation of $\tau$ commutes
with the base change $A \to A_1$ (details omitted; use the
resolution $F_1 \otimes_A A_1 \to F_0 \otimes_A A_1 \to B'_1 \to 0$
to see this). Thus the desired compatibility follows from the result
of the previous paragraph.
\end{proof}

\end{document}
